\section{Objetivos}

\subsection{Objetivo general}
El objetivo general de este trabajo de investigación es desarrollar un flujo de trabajo
procedural e interactivo para la visualización tridimensional de los tensores intermedios en
redes neuronales convolucionales, mediante la integración de Blender y Python, utilizando cuadrículas 
volumétricas jerárquicas dispersas (\textit{OpenVDB}) como alternativa a flujos de trabajo tradicionales
basados en secuencias de imagenes.

\subsection{Objetivos específicos}
Para alcanzar el objetivo general, se plantean los siguientes objetivos específicos:
\begin{itemize}
    \item \textbf{Extraer} y exportar de forma estructurada los mapas de activación y tensores intermedios
    de la arquitectura convolucional YOLOv2 desde su entorno nativo de desarrollo en Python.
    \item \textbf{Transformar} las matrices tridimensionales abstractas en archivos de volumen disperso basados
    en el estándar OpenVDB para optimizar su posterior importación y gestión dentro de Blender.
    \item \textbf{Diseñar} un entorno de visualización procedural interactivo en Blender, utilizando Python y el 
    sistema de nodos de geometría, que permita manipular, animar e inspeccionar la topología original de los datos.
    \item \textbf{Evaluar} la fidelidad y la eficiencia del flujo de trabajo propuesto, demostrando su viabilidad 
    tanto para la interpretación interna del modelo como para la generación de material didáctico ilustrativo.
\end{itemize}